\section{Mean field games, with age as time}\label{sec:mfg}

The Lasry--Lions argument \citep{LasryLions2007} proves uniqueness in a mean field game from two
optimality inequalities and a monotonicity condition on the cost. An earlier survey ruled it out
for the stationary Bewley model, and for a stated reason: the argument cancels a term
$\int(u_1-u_2)\,d\bigl(m_1(0)-m_2(0)\bigr)$ by assuming the two equilibria start from a common
initial distribution, and a stationary discounted equilibrium has no common start --- two
candidate rates carry two different stationary distributions.

Overlapping generations do have one. Every cohort is born with nothing, whatever the rate, so the
date-zero marginal is the same point mass crossed with the invariant earnings law at every
candidate equilibrium. The stated obstruction is simply absent, and the monotonicity condition
becomes a question to check rather than one to dismiss.

\subsection{The game}

Take age as time. A player is a cohort member, a path is
$x=(a_0,\dots,a_{J})$ with $a_0=0$, and the population is a distribution $\eta$ over paths with
date-$t$ marginals $m_t$. The cost at date $t$ is the negative of flow utility, with the budget
substituted in,
\[
  F_t(x,m) \;=\; -\beta^{t}\,u\Bigl(w\bigl(K(m)\bigr)y_t+\bigl(1+r(K(m))\bigr)a_t-a_{t+1}\Bigr),
  \qquad K(m)=\int a\,dm ,
\]
so the population enters through the single scalar $K$, and through it through both factor prices.
This is the price-through-an-aggregate structure of \citet{GraberMatter2024}, except that the
aggregate enters the budget rather than being added to the cost.

\subsection{What the argument actually uses}

Adding the two equilibrium inequalities --- each equilibrium does no worse against its own costs
than the other distribution does --- signs the sum of the date pairings and nothing finer:
\[
  \sum_{t}\;\underbrace{\int\bigl(F_t(\cdot,m_{1t})-F_t(\cdot,m_{2t})\bigr)\,d(m_{1t}-m_{2t})}
  _{\textstyle P_t}\;\le\;0 .
\]

\begin{theorem}[The engine]\label{thm:mfgengine}
Two relaxed equilibria of the finite-horizon game with a common initial distribution have
$\sum_t P_t\le 0$. No monotonicity, no integrability and no continuity is used.
\formal{MeanFieldGame.summedPairing\_nonpos\_of\_pathEquilibrium}
\end{theorem}

The usual hypothesis, that every $F_t$ is Lasry--Lions monotone, is one way to conclude from this
and not the only one: it forces each summand to be zero. The weaker hypothesis that the argument
can carry asks only that the total be nonnegative.

\begin{theorem}[Uniqueness of an aggregate from the summed pairing]\label{thm:mfgagg}
If the cost separates a statistic $A$ of the marginals --- $\sum_tP_t>0$ whenever $A$ differs ---
then two relaxed equilibria from a common initial distribution agree on $A$. No single date is
assumed monotone.
\formal{MeanFieldGame.aggregate\_eq\_of\_pathEquilibrium}
\end{theorem}

Date-by-date monotonicity is the special case in which every summand is already nonnegative
(\lean{summedLasryLionsMonotone\_of\_forall}). The distinction is not cosmetic here, as the next
two subsections show: in this economy the date-by-date hypothesis is false and the summed one is
not.

\subsection{The price channel is monotone; the curvature is not}

Suppose for a moment that $u$ were affine. Then
$F_t(x,m_1)-F_t(x,m_2)=-\beta^t\bigl[(w_1-w_2)y_t+(r_1-r_2)a_t\bigr]$ and the pairing is
\[
  P_t \;=\; -\beta^t (w_1-w_2)\!\int\! y_t\,d(m_{1t}-m_{2t})
        \;-\;\beta^t (r_1-r_2)\!\int\! a_t\,d(m_{1t}-m_{2t}).
\]
The first term is exactly zero: labour is exogenous and the earnings marginal at each age is the
same invariant law in both equilibria. The second is $-\beta^t(r_1-r_2)(A_{1t}-A_{2t})$, and the
factor-price frontier makes $r$ strictly decreasing in $K$, so it is strictly positive at every
date at which the economy with more capital holds more assets. The interest channel is exactly the
Lasry--Lions congestion term: more aggregate capital raises the cost of holding capital. The price
channel of a production economy does not run the wrong way. It is the only part of the problem
that is unambiguously right.

What breaks is the weighting. With $u$ strictly concave the pairing is that same price term
weighted by marginal utility, and the weights do not cancel, because the two economies choose
different savings and so sit at different marginal utilities. Over all pairs of measures the
damage is total: taking concave and convex envelopes in the aggregate shows that unrestricted
monotonicity of this cost forces $u$ to be affine. The question worth asking is therefore the
restricted one --- whether monotonicity holds over the distributions the model can actually
generate.

\subsection{Date by date it fails, and it fails at birth}

Take two economies differing only in $\beta$, each at its own equilibrium, so that each
distribution is priced by the aggregate it actually carries. Since $\beta$ does not appear in
$F_t$, this is a legitimate test of the date-$t$ condition. At the calibration of
Section~\ref{sec:eqm} ($\alpha=9/25$, $\delta=2/25$, $J=60$, Aiyagari's earnings chain), with
$\beta=0.960$ against $\beta=0.950$:

\begin{center}
\begin{tabular}{lrrrrr}
\toprule
$\gamma$ & $0.5$ & $1$ & $2$ & $3$ & $5$\\
\midrule
dates with $P_t<0$ (of $60$) & $3$ & $4$ & $5$ & $6$ & $6$\\
$P_0$ & $-0.00147$ & $-0.00227$ & $-0.00324$ & $-0.00419$ & $-0.00804$\\
$\sum_t P_t$ (undiscounted) & $+0.271$ & $+0.175$ & $+0.091$ & $+0.061$ & $+0.047$\\
\bottomrule
\end{tabular}
\end{center}

The failures are at the start of life, and they survive as the two equilibria close up:
$P_0/(\Delta\beta)^2$ converges to $-25.6$ at $\gamma=1$ and $-44.4$ at $\gamma=3$ as
$\Delta\beta$ falls from $0.02$ to $0.00125$, with the count of negative dates unchanged
throughout. So this is a property of the derivative, not an artefact of a wide gap.

Date zero is the sharp case, and it is worth spelling out, because it is exactly the common
initial distribution that was supposed to rescue the argument. Every cohort is born with nothing,
so $a_0=0$ for all of the mass in both economies, and the affine-utility pairing above is
identically zero: with no assets to price, the aggregate coupling has no grip at all. But the wage
still differs, and the utility gain from the higher wage is weighted by marginal utility, which is
higher in the economy that saves more. At $\gamma=1$ mean date-zero saving is $0.2937$ against
$0.2736$, and $P_0$ is strictly negative. The common start does not make the date-zero pairing
vanish; it removes the one term that could have made it positive.

\subsection{Summed, it holds where the equilibrium lives}

The summed condition is a different matter. Test it over the family the equilibria belong to:
$m(\hat K)$ generated by optimal life-cycle behaviour at the prices of $\hat K$, priced inside $F$
by the aggregate $A(\hat K)$ it actually carries. Every equilibrium is an $m(\hat K)$ with
$A(\hat K)=\hat K$, so monotonicity over pairs from this family is exactly what
Theorem~\ref{thm:mfgagg} needs.

\begin{center}
\begin{tabular}{lrrrr}
\toprule
$\gamma$ & $1$ & $2$ & $3$ & $5$\\
\midrule
pairs with $\sum_t\beta^tP_t<0$, $\hat K/K^*\in[0.5,2]$ & $1/66$ & $0/66$ & $0/66$ & $0/66$\\
pairs failing inside $K\in[3.55,7.40]$ & $0/36$ & $0/36$ & $0/36$ & $0/28$\\
least margin $\sum_t\beta^tP_t/(A_1-A_2)^2$ there & $+0.075$ & $+0.097$ & $+0.121$ & $+0.338$\\
\bottomrule
\end{tabular}
\end{center}

The bracket $K\in[3.55,7.40]$ is the one $r\in[2\%,8\%]$ of Theorem~\ref{thm:exist5}, inside which
the equilibrium is proved to lie. Within it the summed pairing is strictly positive for every pair
and every risk aversion tested, with a margin bounded away from zero, and the margin grows with
$\gamma$. The single failure is at $\gamma=1$ between two candidates one of which sits outside the
bracket entirely.

\subsection{What this leaves}

Three things, in order of how much they are worth.

The negative result is definite and it is not the expected one. Date-by-date Lasry--Lions
monotonicity fails in this economy, for every risk aversion including logarithmic, and the reason
is not that the price channel runs backwards --- it runs the right way, exactly --- but that
strict concavity reweights it, most severely at the ages where there are no assets to price. This
is not a $\mu=1$ frontier of the kind the rest of the programme keeps running into; it fails at
$\mu=1$ too.

The structural gain is real. Theorem~\ref{thm:mfgengine} isolates what the argument does without
any hypothesis, and Theorem~\ref{thm:mfgagg} turns the summed condition directly into uniqueness
of the aggregate, which for this economy is uniqueness outright: one $K$ means one $r$, one
household problem, one answer. The route is no longer dismissed for the wrong reason.

What is missing is a proof of the summed condition from primitives. It is now a measured
conjecture --- one scalar inequality about a whole life, strictly satisfied with room to spare
throughout the bracket in which existence is proved --- rather than a hypothesis known to be
false. That is a better position than the date-by-date condition ever offered, and it is the first
route in this programme whose measured margin improves as risk aversion rises.
